%% kapa: exam=1 round=29 year=2018 period=1 form=A subject=real_estate subject_ko=부동산학원론 seq=3 number=83
%% source: Q-Net 감정평가사 기출문제 (cst00309) artlSeq=5202808 file=extracted/1008/5202808/2018년 제29회 감정평가사 1차시험 기출문제/2018년 제29회 감정평가사 1차시험 기출문제/1교시 A형.pdf
%% conversion: kordoc 4.24.51 -> tools/gen1.py (auto); review=proofread-sonnet; flags=none
\begin{question}{83}
입목에 관한 법령상 옳지 않은 것은?
\begin{choices}
  \item 입목의 소유자는 토지와 분리하여 입목을 양도할 수 있다.
  \item 입목을 위한 법정지상권은 성립하지 않는다.
  \item 토지소유권 또는 지상권 처분의 효력은 입목에 미치지 않는다.
  \item 입목을 목적으로 하는 저당권의 효력은 입목을 베어 낸 경우에 그 토지로부터 분리된 수목에도 미친다.
  \item 지상권자에게 속하는 입목이 저당권의 목적이 되어 있는 경우에는 지상권자는 저당권자의 승낙 없이 그 권리를 포기할 수 없다.
\end{choices}
\end{question}
\answer{2}
